% The box on numerical divergence, \input into LSTNet.tex. The values are from crosstest/authors/replay.json and
% rounding.json and from the scripts in crosstest/divergence (14 October 2026)
\begin{pbox}[breakable]{Same code, same data, same random numbers: why the numbers still part}
Cross-test D gives both codes everything that could make them differ, except the code itself: the artificial data, the
authors' initial weights, their batch order for all 100 epochs (600 steps), and dropout 0, so no random numbers are
drawn at all. Both train in float32: the authors' data, scale and weights are float32 (\fn{torch.zeros}, \fn{.float()},
PyTorch's default), and ours are \fn{Float32} throughout. Still, their numbers part after 10 to 51 epochs:

\begin{valuetable}{@{}lccccc@{}}
\textbf{Case} & \textbf{Epoch kept} & \textbf{Test RSE} & \textbf{Test CORR} & \textbf{Epoch valid RSE} \\
 & Julia / authors & Julia / authors & Julia / authors & \textbf{0.1\,\% apart} \\\midrule
seasonal, $h = 3$  & 100 / 100 & 0.1783 / 0.1800 & 0.9788 / 0.9787 & 14 \\
seasonal, $h = 12$ & 98 / 92   & 0.1909 / 0.1922 & 0.9755 / 0.9750 & 10 \\
trend, $h = 3$     & 99 / 99   & 0.1126 / 0.1094 & 0.8123 / 0.8129 & 51 \\
trend, $h = 12$    & 93 / 79   & 0.1182 / 0.1165 & 0.7517 / 0.7498 & 19 \\
constant, $h = 3$  & 32 / 38   & 0.4426 / 0.4421 & 0.9616 / 0.9618 & 15 \\
constant, $h = 12$ & 98 / 67   & 0.4217 / 0.4174 & 0.9696 / 0.9689 & 21 \\
\end{valuetable}

\nmsub{Not the code: rounding}
Every arithmetic operation in float32 rounds its result to about 7 significant digits. PyTorch and Julia do the same
operations in a different order (sums of many terms, the matrix products of BLAS), so each rounds a little differently,
at every step. After one step the two codes' weights are $2.9\times10^{-8}$ apart (relative distance), after 100 steps
$5.9\times10^{-7}$, after 230 steps $1.8\times10^{-6}$: a drift, still far below any digit printed.

\nmsub{The switch that amplifies it: ReLU}
The convolution's outputs pass through a ReLU, which is a switch: an input just above zero passes, one just below gives
zero, and so does its gradient. When an input lies closer to zero than the gap between the two codes, the codes put it
on opposite sides, and from then on they compute different gradients. We traced the trend case, $h = 3$, step by step:
at every step the 3022 pre-activations that PyTorch had within $10^{-4}$ of zero, compared with Julia's.

\begin{valuetable}{@{}lXl@{}}
\textbf{Step} & \textbf{Event} & \textbf{Weight distance}\\\midrule
1--230 & rounding only & $2.9\times10^{-8}$ to $1.8\times10^{-6}$\\
232 & first flip, filter 5: Julia $+2.6\times10^{-6}$, PyTorch $-1.2\times10^{-7}$ (8 positions) & $1.8\times10^{-6}$\\
249 & flip, filter 5 (5 positions) & $1.8\times10^{-6}$\\
275 & flip, filter 5: Julia $+1.1\times10^{-5}$, PyTorch $-5.2\times10^{-5}$ (10 positions) & $3.7\times10^{-6}$\\
280--300 & more flips, at steps 284, 285, 286, 287, 291, \dots & $8.8\times10^{-6}$ to $1.3\times10^{-4}$\\
400, 600 & & $9.2\times10^{-3}$, $1.6\times10^{-2}$\\
\end{valuetable}

The positions that flip at one step are one stretch of the data seen by several overlapping windows of the batch (data
rows 390, 82 and 517 at steps 232, 249 and 275), not a dead filter: filter 5's weights are as large as the others'.
Why the flip at step 275 starts the growth, and those at 232 and 249 leave the distance where it was, we have not
established.

\nmsub{The authors' code against itself}
How much of this is the translation? We ran the authors' code twice with everything the same but every initial weight
moved by one float32 rounding unit (about $10^{-7}$ relative): up, down, and three times up or down at random. In
lockstep, on the seasonal case, $h = 3$, the two runs are $8\times10^{-8}$ apart at step 1 and $2.6\times10^{-7}$ at
step 115; at step 116 an input $5\times10^{-7}$ from zero flips, and the distance is $1.3\times10^{-6}$, then
$2.1\times10^{-4}$ at step 130, $1.8\times10^{-3}$ at 150 and about $8\times10^{-3}$ at 600. For the trend case the first
flip comes at step 452; for seasonal, $h = 12$, none comes in 600 steps ($6.1\times10^{-7}$ at the end). Over the five
nudged runs, the test RSE of the epoch kept is

\begin{valuetable}{@{}lXl@{}}
\textbf{Case} & \textbf{Test RSE of the authors' code, 5 nudged runs} & \textbf{Epochs kept}\\\midrule
seasonal, $h = 3$  & 0.1800 to 0.1802 & 84 to 100\\
seasonal, $h = 12$ & 0.1922 in all five & 92\\
trend, $h = 3$     & 0.1090 to 0.1108 & 95 to 99\\
trend, $h = 12$    & 0.1165 in all five & 79\\
constant, $h = 3$  & 0.4421 to 0.4422 & 32 or 38\\
constant, $h = 12$ & 0.4172 to 0.4197 & 67 to 85\\
\end{valuetable}

So one rounding unit, once, moves the authors' own results by up to 1.6\,\% (trend, $h = 3$) and the epoch kept by up to
18. Julia against PyTorch differ by rounding at every operation, so they part sooner and in every case, by up to
2.9\,\% (trend, $h = 3$). After parting, neither run gets worse: they settle at different solutions, equally good.

\nmsub{What this means for a replication}
\begin{itemize}
  \item \textbf{Test a translation with the same random numbers, step by step.} Only there can it agree to the last
    digits: here the forecasts to $1.8\times10^{-7}$ (A) and every gradient norm of 3 epochs to $1.1\times10^{-6}$ (B).
  \item \textbf{Measure the original against itself.} A one-unit nudge of its starting point shows how far rounding alone
    moves its results; a translation that differs by no more has nothing left to explain.
  \item \textbf{A single published run carries this spread.} Table~2 reports one run per cell, so its values are uncertain
    by at least as much, before any change of seed, hardware or library.
  \item \textbf{The epoch kept is fragile.} Near the end the validation RSE is nearly flat, so the smallest moves from one
    epoch to another with the least change.
\end{itemize}
The traces are reproducible: \fn{crosstest/divergence/lockstep.py} (the authors' code against itself),
\fn{trace.py} and \fn{trace.jl} (Julia against PyTorch, step by step).
\end{pbox}
